% ===== 3 视觉感知与物理测量 =====
\section{视觉感知与物理测量}
\label{sec:r-visual}

\subsection{成像与尺度标定}

上游视觉模块的目标不是获得“好看的照片”，而是稳定恢复颗粒边界并把像素几何转换为毫米几何。现场采用固定俯视相机、低反光背景和漫射照明，优先降低硬阴影、运动模糊与透视缩放差异。相机高度或焦距改变后必须重新标定尺度；若参考长度 $L_{\mathrm{ref}}$ 在图像中占据 $l_{\mathrm{ref}}$ 像素，则
\begin{equation}
 r=\frac{L_{\mathrm{ref}}}{l_{\mathrm{ref}}}\quad (\mathrm{mm/px}).
\label{eq:r-scale}
\end{equation}
所有轴长由像素量乘以 $r$ 得到。尺度误差近似按比例传递到长短轴和筛分粒径，因此它是场景级系统误差，而不是可以由下游统计自动抵消的噪声。参考物应尽量与测量面共面；若画面不同位置得到的 $r$ 差异明显，则说明单一 mm/px 已不足以描述透视，应先调整相机姿态或加入平面单应/镜头畸变校正。

\subsection{基于 Cellpose-SAM 的实例感知}

自然堆积骨料需要的是\emph{实例分割}而非单纯前景分割：只有每个颗粒对应独立 mask，长短轴、面积和后续质量权重才有明确对象。本文直接复用 ImageGrains 2.0 提供的 Cellpose-SAM 颗粒预训练模型作为感知基座\cite{mair2026imagegrains2,pachitariu2025cellposeSAM}。其优势在于不要求比赛现场重新建立大规模实例标注；输出为整数实例 mask，可继续进行几何测量和人工质检。该预训练模型属于复用的基础能力，本文创新重点不在重新提出分割网络，而在把其实例结果接入物理尺度、筛分等效粒径与质量级配映射。

图\ref{fig:r-demo-detection}给出真实自然堆积样本的实例检测与粒径标注示例。现场质检重点不是逐像素追求视觉完美，而是检查会显著改变级配的结构性错误：整片颗粒被 merge 会产生过大的 $d_i$，单颗粒 split 则制造多个小实例；由于后续质量权重随粒径幂次增长，粗尺度 merge 对最终高分位数尤其敏感。因此正式视觉精度应在人工实例掩膜上同时报告 IoU/F1 与 merge/split，而不能只给一张叠加图。

\begin{figure}[t]
\centering
\includegraphics[width=0.86\linewidth]{agg_001_IG2_full_set_cp_SAM_pred_grains_re_scaled_detection_overview.png}
\caption{真实自然堆积样本的实例检测与粒径标注概览。可视化的主要作用是使最终级配能够回溯到具体实例，并快速发现 merge、split、漏检或明显尺度异常。}
\label{fig:r-demo-detection}
\end{figure}

\subsection{颗粒几何与任务特征}

对实例 mask $M_i$，系统提取面积 $A_i$、周长 $P_i$、凸包面积 $A_{\mathrm{conv},i}$、拟合椭圆长轴 $a_i$ 和短轴 $b_i$。长度经式\eqref{eq:r-scale}转换到毫米，面积按 $r^2$ 转换。面积等效直径定义为
\begin{equation}
 d_{\mathrm{eq},i}=2\sqrt{\frac{A_i}{\pi}},
\label{eq:r-deq}
\end{equation}
并计算三个无量纲形貌量
\begin{equation}
 \mathrm{AR}_i=\frac{b_i}{a_i},\qquad
 C_i=\frac{4\pi A_i}{P_i^2},\qquad
 S_i=\frac{A_i}{A_{\mathrm{conv},i}}.
\label{eq:r-shape}
\end{equation}
其中 AR 描述细长程度，$C$ 描述圆度，$S$ 描述相对于凸包的凹凸程度。这些量一方面参与后续筛分粒径映射，另一方面复用于形貌和异常分支，从而避免为每个任务建立彼此独立、难以解释的模型。

当前投影形貌规则以 AR、$C$ 和 $S$ 将颗粒划分为针片状候选、圆形、棱角状和普通四类；其中“针片状候选”只描述二维投影细长性，不能替代含厚度维度的规范针片状判定。异常分支以物理粒径和凸度做快速筛查：$d>50$~mm 标记为超限，$d<5$~mm 标记为过小候选，40--50~mm 且低凸度的目标仅标记为疑似泥团。前两类有明确几何定义，而泥团本质上是材质语义问题，因此当前几何规则只承担候选筛查，后续若需要可靠泥团识别，应加入颜色/纹理或小型裁剪分类器。

至此，图像已经被压缩为一组可追溯的毫米级颗粒描述 $\{\mathbf{x}_i\}$。但这些几何量仍然不是标准筛分结果：短轴是二维投影量，机械筛孔对应三维动态通过行为；同时图像按实例计数，而机械筛分按质量称重。下一节处理全文最关键的“筛分语义鸿沟”。
